\chapter{Eksperimen dan Analisis}

Bab ini memaparkan seluruh tahapan eksperimen yang dilakukan, mulai dari perancangan rentang kontrol,
identifikasi sistem, hingga evaluasi perbandingan performa kontroler dalam
skenario \emph{closed-loop}.
Setiap keputusan perancangan disertai justifikasi kuantitatif dari data eksperimen.

% =============================================================================
% DATA COMMANDS — ganti nilai di sini jika data eksperimen berubah
% =============================================================================

% --- Open-loop ---
\newcommand{\openloopsamples}{8146}
\newcommand{\openloopsteps}{122}
\newcommand{\openloopinterval}{1}

% --- Capacity benchmark ---
\newcommand{\tfetchms}{203}
\newcommand{\batchminCap}{3}
\newcommand{\batchmaxCap}{2036}
\newcommand{\pollminCap}{604}
\newcommand{\pollmaxCap}{1599}
\newcommand{\kneepoint}{500}
\newcommand{\kneeCpu}{31.5}
\newcommand{\docsizeKB}{0.463}

% --- Resource config ---
\newcommand{\escpus}{1,0}
\newcommand{\esmemory}{1024 MB}
\newcommand{\esjavaheap}{512 MB}
\newcommand{\esversion}{8.10.0}

% --- Closed-loop run ---
\newcommand{\closedlooprun}{closed\_loop\_20260415\_231610}
\newcommand{\closedloopsamples}{1930}

% --- Normalization override (LQR) ---
\newcommand{\normQueueMean}{0}
\newcommand{\normQueueStd}{5000}
\newcommand{\normBatchMean}{250}
\newcommand{\normBatchStd}{250}
\newcommand{\normInvPollMean}{5,0}
\newcommand{\normInvPollStd}{2,0}

% =============================================================================

\section{Tujuan Eksperimen}

Eksperimen pada penelitian ini bertujuan untuk menjawab empat pertanyaan utama.
Pertama, bagaimana rentang operasional yang tepat untuk parameter kontrol berdasarkan kapasitas sistem
yang terukur, bukan berdasarkan heuristik intuitif.
Kedua, apakah dinamika sinkronisasi CQRS dapat dimodelkan cukup baik untuk mendukung rancangan
kontroler berbasis model.
Ketiga, seberapa efektif LQR dibandingkan kontroler lain (berbasis aturan dan PID) dalam
menjaga antrean tetap rendah pada berbagai pola beban.
Keempat, bagaimana setiap kontroler berperilaku pada skenario beban yang berbeda dan apa yang
membedakan konsistensi masing-masing.

\section{Lingkungan Eksperimen}

Seluruh eksperimen dijalankan pada lingkungan yang sama untuk menjaga konsistensi hasil.
Arsitektur sistem uji mengikuti rancangan pada Bab~III, terdiri dari \emph{write model} (PostgreSQL),
\emph{message broker} (Kafka dengan Debezium CDC), \emph{message sink}, dan \emph{read database}
(Elasticsearch).

Komponen \emph{read database} menggunakan Elasticsearch versi~\esversion{} dengan alokasi~\escpus{}
CPU, memori~\esmemory{}, dan Java heap~\esjavaheap{} sebagai konfigurasi dasar.
Pada eksperimen \emph{closed-loop}, konfigurasi sumber daya divariasikan menjadi dua tingkat
utama (penuh dan setengah) pada dua mesin uji berbeda untuk menguji portabilitas kontroler
terhadap perubahan kapasitas sistem dan perangkat keras (lihat Subbab~\ref{sec:hasil-full},
\ref{sec:hasil-half}, dan~\ref{sec:robustness}).
Selain itu, konfigurasi seperempat diuji sebagai kasus batas yang dibahas singkat pada
Subbab~\ref{sec:quarter}.

\section{Penentuan Rentang Kontrol dari Kapasitas Sistem}

Sebelum menjalankan eksperimen identifikasi sistem, terlebih dahulu dilakukan pengukuran kapasitas
\emph{throughput} untuk menentukan rentang operasional yang terjustifikasi secara empiris.

\subsection{Pengukuran Kapasitas (\emph{Capacity Benchmark})}

Eksperimen \emph{capacity benchmark} dilakukan dengan memvariasikan \emph{batch size} pada tujuh level (1, 10, 50, 100, 250, 500, 1000) dengan
\emph{poll interval} tetap 500~ms.
Untuk setiap level, \emph{throughput} aktual (msg/s), utilisasi CPU, dan waktu pengindeksan per
dokumen dicatat selama 30 detik dalam kondisi antrean selalu tersedia.

Hasil pengukuran menunjukkan bahwa ukuran dokumen rata-rata adalah~\docsizeKB~KB dan waktu
\emph{round-trip} Kafka (\emph{t\_fetch}) adalah~\tfetchms~ms.
Nilai~\emph{t\_fetch} ini menjadi parameter kunci dalam perancangan rentang \emph{poll interval}
(lihat Subbab~\ref{sec:poll-range}).

\begin{table}[H]
    \centering
    \caption{Hasil pengukuran \emph{throughput} dan utilisasi CPU per level \emph{batch size}}
    \label{tab:capacity-benchmark}
    \renewcommand{\arraystretch}{1.4}
    \begin{tabular}{|r|r|r|r|}
    \hline
    \textbf{Batch} & \textbf{Throughput (msg/s)} & \textbf{CPU (\%)} & \textbf{t\_index (ms/dok)} \\
    \hline
    1     & 56      & 14,6 & 30,97 \\
    \hline
    10    & 590     & 16,5 &  3,02 \\
    \hline
    50    & 2.499   & 26,0 &  0,77 \\
    \hline
    100   & 4.582   & 28,7 &  0,39 \\
    \hline
    250   & 8.788   & 28,2 &  0,27 \\
    \hline
    \textbf{500}   & \textbf{14.132}  & \textbf{31,5} &  \textbf{0,16} \\
    \hline
    1000  & 15.378  & 63,7 &  0,26 \\
    \hline
    \end{tabular}
\end{table}

Dari Tabel~\ref{tab:capacity-benchmark} terlihat bahwa \emph{throughput} meningkat nyata hingga
\emph{batch}~=~\kneepoint, dengan CPU~\kneeCpu\%.
Setelah titik ini, kenaikan \emph{throughput} hanya 9\% (dari 14.132 ke 15.378~msg/s) namun CPU
melonjak dua kali lipat (dari 32\% ke 64\%).
Titik ini disebut \emph{knee point} dan menjadi batas atas yang efektif untuk operasi yang efisien.

Gambar~\ref{fig:knee-point} memperlihatkan pola ini secara visual.
Kurva \emph{throughput} (biru) tumbuh hampir linier hingga \emph{batch} = 500, lalu hampir
mendatar setelahnya, sementara kurva CPU (merah, sumbu kanan) baru melonjak tajam setelah
melewati titik tersebut.

\begin{figure}[H]
    \centering
    \includegraphics[width=0.85\textwidth]{IV_R10__knee_point.png}
    \caption{Identifikasi \emph{knee point} dari hasil \emph{capacity benchmark}}
    \label{fig:knee-point}
\end{figure}

\subsection{Penentuan Rentang \emph{Batch Size}}
\label{sec:batch-range}

Rentang \emph{batch size} diturunkan dari dua batasan fisik menggunakan Bryson's rule.
Bryson's rule adalah kaidah perancangan dalam kontrol optimal yang menetapkan nilai bobot
atau batasan berdasarkan rentang operasional yang dapat diterima, bukan dari penalaan
heuristik (Bryson dan Ho, 1975).
Pada konteks ini, kaidah tersebut diterapkan untuk menentukan batas atas \emph{batch size}
agar parameter kontrol selalu berada dalam wilayah operasi yang aman.
Setiap batasan memberikan jumlah dokumen maksimum per \emph{batch} yang masih aman, dan yang
mengikat adalah yang paling kecil di antara keduanya.

\begin{enumerate}
\item \textbf{Batasan memori:}
$\text{batch\_max}_{\text{mem}} = \lfloor J \times \alpha / d \rfloor$,
dengan $J = 512$~MB (\emph{heap}), $\alpha = 0{,}10$ (fraksi heap untuk \emph{bulk buffer}),
dan $d = \text{\docsizeKB}$~KB (ukuran dokumen terukur).
Hasilnya $\text{batch\_max}_{\text{mem}} = 113\,197$ dokumen.

\item \textbf{Batasan CPU:}
$\text{batch\_max}_{\text{cpu}} = C_{\text{limit}} \times 1000 / t_{\text{index}}$,
dengan $C_{\text{limit}} = 1{,}0$~core dan $t_{\text{index}} = 0{,}393$~ms/dok, yaitu waktu
rata-rata yang dibutuhkan Elasticsearch untuk mengindeks satu dokumen pada kondisi \emph{bulk},
diukur dari kolom $t\_\text{index}$ pada Tabel~\ref{tab:capacity-benchmark}.
Hasilnya $\text{batch\_max}_{\text{cpu}} = 2\,545$ dokumen.

\item \textbf{Batasan efektif:}
$\text{batch\_max} = \min(113\,197,\; 2\,545) \times 0{,}80 = \batchmaxCap$.
Karena CPU jenuh jauh lebih cepat daripada memori (2\,545 jauh lebih kecil dari 113\,197),
CPU menjadi batasan yang mengikat.
\emph{Batch} di atas 2\,545 hanya akan menambah waktu pemrosesan tanpa benar-benar menambah
\emph{throughput}, dan justru berisiko menyebabkan antrean penuh di tahap berikutnya.
Faktor keamanan 0,80 menyisakan ruang untuk \emph{garbage collection} dan proses \emph{merge}.

\item \textbf{Batas bawah:}
$\text{batch\_min} = \lceil t_{\text{conn}} / t_{\text{index}} \rceil = \lceil 1/0{,}39 \rceil = \batchminCap$,
dengan $t_{\text{conn}} = 1$~ms adalah estimasi biaya tetap pengiriman satu permintaan
\emph{bulk} HTTP pada koneksi \emph{keep-alive} (tidak termasuk TCP \emph{handshake}
yang sudah terjadi sebelumnya).
Nilai ini konsisten dengan observasi bahwa pada \emph{batch} = 1, overhead pengiriman
mendominasi waktu pengindeksan aktual, sehingga mengirim setidaknya 3 dokumen per
permintaan sudah cukup untuk menyeimbangkan biaya koneksi.
Nilai ini adalah titik di mana \emph{overhead} koneksi melebihi manfaat pemrosesan \emph{batch}.
\end{enumerate}

Rentang \emph{batch size} yang terjustifikasi adalah $[\batchminCap,\; \batchmaxCap]$.

\subsection{Penentuan Rentang \emph{Poll Interval}}
\label{sec:poll-range}

Rentang \emph{poll interval} ditentukan dari dua arah.

\begin{enumerate}
\item \textbf{Batas bawah (floor):}
$\text{poll\_min} = t_{\text{fetch}} + \text{batch\_mid} \times t_{\text{index}} + t_{\text{overhead}}$,
dengan $t_{\text{fetch}} = \text{\tfetchms}$~ms (waktu \emph{round-trip} Kafka yang terukur),
$\text{batch\_mid} = (\batchmaxCap + \batchminCap)/2 \approx 1019$, $t_{\text{index}} = 0{,}393$~ms/dok,
dan $t_{\text{overhead}} = 8$~ms adalah total biaya tetap per siklus \emph{poll} di luar waktu
\emph{fetch} dan pengindeksan, mencakup overhead \emph{consumer} Kafka, pengiriman \emph{bulk} HTTP,
dan \emph{parsing} respons.
Hasilnya: $\text{poll\_min} = 203 + 1019 \times 0{,}393 + 8 \approx \pollminCap$~ms.
Di bawah nilai ini, \emph{poll} selesai sebelum data dari Kafka siap, sehingga hanya membuang sumber daya.

\item \textbf{Batas atas (ceiling):}
$\text{poll\_max} = L_{\text{SLA}} - \text{batch\_mid} \times t_{\text{index}}$,
dengan $L_{\text{SLA}} = 2000$~ms (latensi end-to-end yang dapat diterima).
Hasilnya: $\text{poll\_max} = 2000 - 401 \approx \pollmaxCap$~ms.
Di atas nilai ini, latensi sinkronisasi melebihi batas yang dapat diterima.
\end{enumerate}

Rentang \emph{poll interval} yang terjustifikasi adalah $[\pollminCap,\; \pollmaxCap]$~ms.

\subsection{Penetapan Titik Operasional Nominal ($u_{\text{nom}}$)}

Titik nominal $u_{\text{nom}} = [\text{batch\_nom},\; \text{inv\_poll\_nom}]$ dipilih sebagai
titik keseimbangan yang memberikan kapasitas cukup tanpa mendekati batas saturasi CPU.

\begin{enumerate}
\item \textbf{batch\_nom = 250.}
Dari hasil \emph{benchmark}, \emph{batch}~=~250 menghasilkan \emph{throughput} 8.788~msg/s dengan
CPU 28\%, yakni di bawah \emph{knee point} (batch~=~\kneepoint, CPU~=~\kneeCpu\%).
Titik ini memiliki ruang untuk naik menuju \emph{knee point} jika antrean menumpuk,
sekaligus ruang untuk turun jika sumber daya tertekan.

\item \textbf{inv\_poll\_nom = 5{,}0 (poll = 200~ms).}
\emph{Poll interval} 200~ms mendekati $t_{\text{fetch}} = \text{\tfetchms}$~ms.
Pada \emph{poll} $\geq t_{\text{fetch}}$, setiap siklus \emph{poll} memiliki peluang tinggi menemukan
data baru (efisiensi \emph{polling} $= \min(1,\; t_{\text{fetch}} / \text{poll}) = 1{,}0$).
\emph{Poll} lebih cepat dari $t_{\text{fetch}}$ hanya akan membuang siklus CPU.
\end{enumerate}

\section{Eksperimen Open-Loop dan Identifikasi Sistem}

\subsection{Prosedur Eksperimen}

Eksperimen \emph{open-loop} dilakukan menggunakan grid 10$\times$10 dengan
spasi logaritmik dalam rentang yang telah ditetapkan.
Nilai \emph{batch size} yang diuji adalah \{3, 6, 13, 26, 54, 112, 232, 478, 987, 2036\} dan
\emph{poll interval} \{604, 673, 750, 836, 931, 1037, 1156, 1288, 1435, 1599\}~ms.
Spasi logaritmik dipilih karena respons sistem (terutama \emph{throughput}) menuju saturasi
terhadap kenaikan \emph{batch size}.

Setiap kombinasi kontrol diaplikasikan selama 60 detik setelah pra-pengisian antrean,
kemudian diikuti periode \emph{settling} 10 detik.
Total menghasilkan \openloopsamples{} sampel dari \openloopsteps{} langkah.

Variabel state yang dicatat adalah panjang antrean (\emph{queue\_length}),
utilisasi CPU, utilisasi memori, dan operasi I/O tulis.
Latensi dicatat sebagai variabel observasi saja dan tidak dimasukkan ke dalam vektor state,
karena pengujian sebelumnya menunjukkan bahwa latency di vektor state menyebabkan LQR menekan
\emph{batch size} secara berlebihan (\emph{batch collapse}).

Variabel kontrol yang divariasikan adalah \emph{batch size} dan \emph{inv\_poll\_interval}
($= 1000 / \text{poll\_ms}$).
Penggunaan \emph{inv\_poll\_interval} sebagai variabel kontrol (bukan \emph{poll\_interval})
mempertahankan hubungan yang lebih linier antara variabel ini dan \emph{throughput},
karena \emph{throughput} proporsional dengan frekuensi \emph{polling}, bukan intervalnya.

\subsection{Kualitas Model Identifikasi}

Model state-space $\mathbf{x}_{k+1} = A\mathbf{x}_k + B\mathbf{u}_k$ diestimasi menggunakan
regresi kuadrat terkecil (\emph{least squares}).
Sebelum estimasi, data difilter untuk membuang fase \emph{settle} (transien antar
konfigurasi) dan sampel dengan antrean nol (\emph{queue exhausted}).
Data tersisa dibagi menjadi 80\% sebagai data training dan 20\% sebagai data uji, dengan
pembagian dilakukan per kombinasi parameter \emph{(batch, poll)} untuk mencegah kebocoran
temporal antar himpunan.
Matriks $A$ dan $B$ kemudian diestimasi dari data training, dan kualitas fit ($R^2$) dihitung
pada data tersebut (\emph{in-sample fit}).
Nilai $R^2$ merepresentasikan seberapa baik model linier $A\mathbf{x}_k + B\mathbf{u}_k$
menjelaskan transisi state pada data training.

\begin{table}[H]
    \centering
    \caption{Kualitas fit model state-space per variabel}
    \label{tab:sysid-fit}
    \renewcommand{\arraystretch}{1.4}
    \begin{tabular}{|l|c|l|}
    \hline
    \textbf{Variabel State} & \textbf{$R^2$} & \textbf{Interpretasi} \\
    \hline
    Queue length & 0,994 & Sangat baik (mendekati integrator sempurna) \\
    \hline
    Utilisasi CPU & 0,694 & Cukup \\
    \hline
    Operasi I/O & 0,206 & Rendah \\
    \hline
    Utilisasi memori & 0,316 & Rendah (dinamika didominasi siklus GC JVM, kurang responsif terhadap perubahan batch size dan poll interval) \\
    \hline
    \end{tabular}
\end{table}

$R^2$ memori yang rendah konsisten dengan fakta bahwa heap JVM dikelola secara internal oleh
\emph{garbage collector} dan hampir tidak terpengaruh oleh \emph{batch size} atau
\emph{poll interval}.
Meski demikian, memori tetap dimasukkan ke dalam vektor state dengan bobot kecil pada matriks
$Q$ agar deviasi yang ekstrem tetap diperhitungkan dalam fungsi biaya, walaupun model linier
tidak dapat memprediksi dinamikanya dengan akurat.

\subsection{Koreksi Matriks B dari Data Kapasitas}

Matriks B yang diperoleh dari eksperimen \emph{open-loop} dengan eksitasi acak menunjukkan masalah
pada baris \emph{queue}: $B[\text{queue}, \text{batch}]$ bernilai positif pada estimasi awal,
yang secara fisik tidak masuk akal (batch besar seharusnya menguras antrean, bukan memperbesarnya).
Penyebabnya adalah \emph{confounding}: \emph{batch size} berkorelasi dengan besarnya pra-pengisian
antrean dalam prosedur eksperimen, sehingga regresi linier menangkap hubungan yang tidak kausal.

Untuk mengatasi hal ini, baris $B[\text{queue}, :]$ dikoreksi menggunakan data \emph{capacity benchmark}:

\begin{equation}
B[\text{queue}, \text{batch}] = -\frac{(\text{tput}(\mu_b + \sigma_b) - \text{tput}(\mu_b)) \cdot \Delta t}{\sigma_q}
\end{equation}

\begin{equation}
B[\text{queue}, \text{inv\_poll}] = -\frac{\text{tput}(\mu_b) \cdot (\text{eff}(\mu_p + \sigma_p) - \text{eff}(\mu_p)) \cdot \Delta t}{\sigma_q}
\end{equation}

dengan $\text{tput}(\cdot)$ adalah fungsi \emph{throughput} yang diinterpolasi dari data benchmark,
$\text{eff}(\cdot) = \min(1,\; t_{\text{fetch}} / \text{poll\_ms})$ adalah efisiensi \emph{polling},
$\mu_b = 130$, $\sigma_b = 75$ adalah mean dan standar deviasi \emph{batch} dari sysid,
$\mu_p = 5{,}0$, $\sigma_p = 3{,}0$ adalah mean dan standar deviasi \emph{inv\_poll},
$\sigma_q = 10.000$ adalah standar deviasi \emph{queue}, dan $\Delta t = 1$~detik.

Nilai yang dihasilkan adalah $B[\text{queue}, \text{batch}] = -0{,}21$ dan
$B[\text{queue}, \text{inv\_poll}] \approx 0$ (karena pada titik nominal \emph{poll} $\approx t_{\text{fetch}}$,
efisiensi sudah jenuh pada 1,0 sehingga penambahan \emph{inv\_poll} tidak meningkatkan \emph{throughput}).

Implikasi dari $B[\text{queue}, \text{inv\_poll}] \approx 0$ bagi DARE adalah bahwa kanal
\emph{inv\_poll} tidak berkontribusi pada drainase antrean, sehingga seluruh upaya koreksi
antrean dialokasikan ke kanal \emph{batch size}.
Ini menghasilkan gain $K$ yang lebih agresif pada \emph{batch}, sesuai dengan perilaku yang diinginkan.

\section{Rancangan Matriks Biaya LQR}

\subsection{Matriks R (Penalti Usaha Kontrol)}

Matriks $R$ dirancang menggunakan \emph{Bryson's rule} berbasis rentang operasional:

\begin{equation}
R_{ii} = \frac{R_{\text{base}}}{(\Delta u_{i,\max,\text{norm}})^2}
\end{equation}

dengan $\Delta u_{i,\max} = \max(u_{i,\max} - u_{i,\text{nom}},\; u_{i,\text{nom}} - u_{i,\min})$
adalah maksimum simpangan dari nominal, dan $\Delta u_{i,\max,\text{norm}} = \Delta u_{i,\max} / \sigma_{u_i}$
adalah simpangan tersebut dalam satuan normalisasi.

Dengan rentang dari \emph{capacity benchmark}:
$\Delta u_{\text{batch},\max} = \max(2036 - 250,\; 250 - 3) = 1786$ dan
$\Delta u_{\text{inv\_poll},\max} = \max(10 - 5,\; 5 - 0{,}625) = 5$,
serta $\sigma_{\text{batch}} = 75$ dan $\sigma_{\text{inv\_poll}} = 3$,
diperoleh:

\begin{equation}
R = R_{\text{base}} \cdot \text{diag}\!\left(\frac{1}{(1786/75)^2},\; \frac{1}{(5/3)^2}\right)
= 1{,}0 \cdot \text{diag}(1{,}76 \times 10^{-3},\; 0{,}36)
\end{equation}

Nilai $R[\text{batch}] \ll R[\text{inv\_poll}]$ mencerminkan fakta bahwa rentang operasional
\emph{batch size} ($\sim 2000$) jauh lebih besar daripada rentang \emph{inv\_poll} ($\sim 5$),
sehingga perubahan \emph{batch} ``lebih murah'' secara proporsional dan DARE dapat menggunakannya
lebih leluasa.

\subsection{Matriks Q (Penalti Deviasi State)}

Tiga preset matriks Q dirancang menggunakan \emph{Bryson's rule} pada ruang ternormalisasi,
untuk skenario prioritas yang berbeda.
Latensi tidak dimasukkan sebagai state (hanya sebagai variabel observasi), sehingga vektor state
adalah $\mathbf{x} = [\text{queue, cpu, mem, io}]^T$.
Diagonal Q dinyatakan sebagai bobot tanpa dimensi karena perhitungan biaya dilakukan pada ruang
$z$-score, sehingga rasio antar bobot mencerminkan prioritas relatif.

Setiap elemen diagonal $Q$ diturunkan dari Bryson's rule pada ruang ternormalisasi sebagai:

\begin{equation}
Q_{ii} = \left( \frac{\sigma_i}{\Delta x_{i,\text{maks}}} \right)^2 \times m_i
\label{eq:q-bryson}
\end{equation}

dengan $\sigma_i$ adalah standar deviasi ruang operasional (lihat Tabel~\ref{tab:lqr-norm}),
$\Delta x_{i,\text{maks}}$ adalah deviasi maksimum yang dapat diterima sistem, dan $m_i$ adalah
\emph{multiplier} yang menyatakan prioritas relatif antar preset.
Pemilihan $\Delta x_{i,\text{maks}}$ untuk masing-masing variabel state ditampilkan pada
Tabel~\ref{tab:q-input}.

\begin{table}[H]
    \centering
    \caption{Parameter input penurunan $Q_\text{base}$ menggunakan Bryson's rule}
    \label{tab:q-input}
    \renewcommand{\arraystretch}{1.4}
    \begin{tabular}{|l|r|r|p{6cm}|}
    \hline
    \textbf{Variabel} & \textbf{$\sigma_i$} & \textbf{$\Delta x_{i,\text{maks}}$} & \textbf{Justifikasi $\Delta x_\text{maks}$} \\
    \hline
    queue\_length & 5000 & 5000 & Deviasi sebesar satu standar deviasi sudah dianggap batas yang ketat untuk antrean \\
    \hline
    cpu\_util & 25 & 30 & Sedikit di atas standar deviasi, mencerminkan toleransi lonjakan CPU singkat \\
    \hline
    mem\_util & 20 & 100 & Memori kurang responsif terhadap aksi kontrol sehingga deviasi besar masih dapat ditoleransi \\
    \hline
    io\_write\_ops & 40 & 50 & Sekitar 1{,}25 standar deviasi, sesuai rentang p95 yang teramati \\
    \hline
    \end{tabular}
\end{table}

Substitusi nilai-nilai tersebut menghasilkan vektor $Q_\text{base}$:

\begin{align*}
Q_\text{base, queue} &= (5000/5000)^2 = 1{,}000 \\
Q_\text{base, cpu}   &= (25/30)^2     = 0{,}694 \\
Q_\text{base, mem}   &= (20/100)^2    = 0{,}040 \\
Q_\text{base, io}    &= (40/50)^2     = 0{,}640
\end{align*}

Tiga preset diperoleh dengan mengalikan $Q_\text{base}$ dengan vektor \emph{multiplier} $m$
yang menggeser prioritas antar variabel.
Nilai 4 dipilih karena memberikan perbedaan bobot yang cukup besar untuk menghasilkan perilaku
kontrol yang berbeda secara nyata, namun tetap dalam orde besaran yang sama agar matriks $K$
yang dihitung tetap stabil secara numerik.
Tabel~\ref{tab:lqr-q-presets} merangkum kombinasi tersebut.

\begin{table}[H]
    \centering
    \caption{Preset matriks $Q$ sebagai $Q_\text{base} \times m$ dan prioritas relatif}
    \label{tab:lqr-q-presets}
    \renewcommand{\arraystretch}{1.4}
    \begin{tabular}{|c|c|l|c|}
    \hline
    \textbf{Preset} & \textbf{Multiplier $m$} & \textbf{Diagonal Q} & \textbf{Prioritas relatif} \\
    \hline
    Q1 & [4, 1, 1, 1] & diag(4{,}000, 0{,}694, 0{,}040, 0{,}640) & 74\%, 13\%, 1 \%, 12 \% \\
    \hline
    Q2 & [1, 4, 4, 4] & diag(1{,}000, 2{,}778, 0{,}160, 2{,}560) & 15\%, 43\%, 2 \%, 39 \% \\
    \hline
    Q3 & [4, 4, 4, 4] & diag(4{,}000, 2{,}778, 0{,}160, 2{,}560) & 42\%, 29\%, 2 \%, 27 \% \\
    \hline
    \end{tabular}
\end{table}

Q1 memprioritaskan antrean sehingga kontroler menguras antrean lebih cepat dengan
konsekuensi pemakaian sumber daya yang lebih agresif.
Q2 membobot tinggi CPU dan IO, sehingga utilisasi sumber daya tetap rendah meskipun antrean
turun lebih lambat.
Q3 memperhatikan antrean dan sumber daya secara setara, cocok ketika kecepatan sinkronisasi
penting tanpa mengorbankan stabilitas sumber daya.

Bobot memori tetap kecil di seluruh preset karena memori dikelola oleh \emph{garbage collector}
JVM dan hampir tidak dipengaruhi secara langsung oleh \emph{batch size} atau \emph{poll interval}.
Memasukkan bobot besar pada memori hanya akan menambah \emph{noise} pada sinyal kontrol.

\subsection{State Target dan Titik Nominal Kontrol}

State target $\mathbf{x}^*$ dipilih sebagai berikut:

\begin{table}[H]
    \centering
    \caption{State target dan justifikasi pemilihan}
    \label{tab:state-target}
    \renewcommand{\arraystretch}{1.4}
    \begin{tabular}{|l|r|l|}
    \hline
    \textbf{Variabel} & \textbf{Target} & \textbf{Justifikasi} \\
    \hline
    queue\_length & 0 & Kondisi ideal: tidak ada antrean tertunda \\
    \hline
    cpu\_util & 5\% & Mendekati \emph{baseline} CPU saat idle (\emph{one-sided}: hanya dihukum jika $>5\%$) \\
    \hline
    mem\_util & 50\% & Nilai tengah rentang operasional; \emph{one-sided} \\
    \hline
    io\_write\_ops & 30 ops/s & Nilai median dari data historis; \emph{one-sided} \\
    \hline
    \end{tabular}
\end{table}

Penalti \emph{one-sided} pada cpu, mem, dan io berarti variabel ini hanya dihukum jika nilainya
melebihi target (bukan juga ketika lebih rendah).
Hal ini mencerminkan realitas bahwa CPU rendah bukan masalah, melainkan CPU tinggi yang perlu dihindari.

\subsection{Normalisasi untuk Ruang Kerja LQR}
\label{sec:lqr-norm}

Saat eksperimen \emph{open-loop} dijalankan dengan antrean besar (0--200.000), standar deviasi
$\sigma_q = 10.000$ mencerminkan rentang yang sangat lebar.
Namun dalam operasi \emph{closed-loop} aktual, antrean hanya bergerak dalam kisaran 0--7.000
(persentil ke-95 = 7.049).

Jika normalisasi sysid digunakan apa adanya, maka kesalahan state $\Delta q = 5.000$ dalam ruang
normalisasi hanya bernilai $5.000/10.000 = 0{,}5$, sehingga aksi kontrol yang dihasilkan DARE
jauh lebih kecil dari yang dibutuhkan.
Akibatnya, LQR terlihat tidak responsif meskipun gain $K$-nya besar.

Untuk mengatasi hal ini, normalisasi LQR di-\emph{override} ke rentang operasional
\emph{closed-loop} yang aktual.
Pada tabel berikut, kolom \textbf{Target ($x^*$)} adalah titik referensi yang dijadikan pusat
normalisasi, bukan rata-rata empiris dari data.
Misalnya, target queue\_length adalah 0 karena kontroler ingin menggiring antrean ke nol,
bukan ke nilai rata-rata historisnya.
Standar deviasi tetap diturunkan dari rentang operasional aktual.

\begin{table}[H]
    \centering
    \caption{Normalisasi LQR (di-\emph{override} dari sysid ke rentang \emph{closed-loop})}
    \label{tab:lqr-norm}
    \renewcommand{\arraystretch}{1.4}
    \begin{tabular}{|l|r|r|l|}
    \hline
    \textbf{Variabel} & \textbf{Target ($x^*$)} & \textbf{Std} & \textbf{Justifikasi std} \\
    \hline
    queue\_length & \normQueueMean & \normQueueStd & $\text{p95}_{cl} / 1{,}645 \approx 4286$ dibulatkan ke 5000 \\
    \hline
    cpu\_util & 5,0 & 25,0 & $\text{p95}_{cl}(=50\%) / 2 = 25$ (mencakup 5--55\%) \\
    \hline
    mem\_util & 50,0 & 20,0 & Sama dengan sysid (p5--p95 $\approx$ 21--73\%) \\
    \hline
    io\_write\_ops & 30,0 & 40,0 & $\text{p95}_{cl}(=99) / 2{,}5 = 40$ (mencakup 0--100) \\
    \hline
    batch\_size & \normBatchMean & \normBatchStd & $\text{max\_useful}/4 = 2036/4 \approx 509$ dibulatkan ke 250 \\
    \hline
    inv\_poll\_interval & \normInvPollMean & \normInvPollStd & Std aktual CL = 1,7; dibulatkan ke 2,0 \\
    \hline
    \end{tabular}
\end{table}

Dengan normalisasi ini, kesalahan antrean 5.000 dalam ruang normalisasi bernilai
$5.000 / 5.000 = 1{,}0$ (satu standar deviasi), menghasilkan aksi koreksi \emph{batch} yang
proporsional dengan kebutuhan aktual.

\section{Hasil Eksperimen Closed-Loop}

\subsection{Setup Eksperimen}

Eksperimen \emph{closed-loop} dijalankan pada satu mesin uji berbasis prosesor AMD Ryzen
dengan seluruh komponen sistem dijalankan sebagai kontainer Docker.
Tumpukan komponen meliputi PostgreSQL sebagai \emph{write database}, Kafka dengan Debezium CDC
sebagai \emph{message broker}, sebuah \emph{message sink} sebagai pengkonsumsi event, dan
Elasticsearch sebagai \emph{read database}.

Untuk menguji robustness kontroler terhadap perubahan kapasitas sistem, eksperimen dijalankan
pada tiga variasi sumber daya Elasticsearch.
Konfigurasi penuh (\emph{Full}) mengalokasikan 2{,}0 CPU, 2048 MiB memori, dan JVM heap 1024 MB.
Konfigurasi setengah (\emph{Half}) menggunakan separuh dari semuanya, yaitu 1{,}0 CPU,
1024 MiB memori, dan JVM heap 512 MB.
Konfigurasi seperempat (\emph{Quarter}) digunakan sebagai kasus batas dengan 0{,}5 CPU,
512 MiB memori, dan JVM heap 256 MB.

Pada masing-masing konfigurasi diuji enam strategi pengendalian, yaitu Statik, Rule-Based,
PID, LQR-Q1, LQR-Q2, dan LQR-Q3.
Setiap kombinasi kontroler dan pola beban dijalankan satu kali pada lima pola, yaitu
\emph{step}, \emph{ramp}, \emph{impulse}, \emph{periodic\_step}, dan \emph{step\_low},
sehingga total terdapat 30 percobaan per konfigurasi sumber daya.

\subsection{Metrik Evaluasi: Fungsi Biaya $J$ dan Regret Ternormalisasi}

Perbandingan utama menggunakan fungsi biaya $J$ yang dihitung secara \emph{post-hoc} pada
seluruh sampel dalam jendela waktu yang ditetapkan.

\begin{equation}
J = \frac{1}{T} \sum_{k=1}^{T} \left[(\mathbf{x}_k - \mathbf{x}^*)^\top Q (\mathbf{x}_k - \mathbf{x}^*) + (\mathbf{u}_k - \mathbf{u}_{\text{nom}})^\top R (\mathbf{u}_k - \mathbf{u}_{\text{nom}})\right]
\end{equation}

Nilai $J$ yang lebih rendah berarti kontroler lebih dekat ke kondisi yang diinginkan.
\emph{Regret} ternormalisasi dihitung sebagai
$r = (J_{\text{controller}} - J_{\text{best}}) / J_{\text{best}}$,
dengan $J_{\text{best}}$ adalah nilai $J$ kontroler terbaik pada konfigurasi yang sama.
Nilai $r = 0$ berarti kontroler terbaik dan $r > 0$ berarti lebih tinggi dari terbaik
secara proporsional.

\subsection{Hasil Sumber Daya Penuh}
\label{sec:hasil-full}

Pada konfigurasi penuh, sistem mengonsumsi 1724 sampel dengan latensi rata-rata 1848~ms
di seluruh kontroler dan pola beban.
Tabel~\ref{tab:cost-J-full} menunjukkan rangkuman fungsi biaya $J$ rata-rata per kontroler
pada masing-masing preset $Q$, beserta \emph{regret} ternormalisasi rata-rata.

\begin{table}[H]
    \centering
    \caption{Nilai rata-rata fungsi biaya $J$ pada konfigurasi sumber daya penuh}
    \label{tab:cost-J-full}
    \renewcommand{\arraystretch}{1.4}
    \begin{tabular}{|l|r|r|r|r|}
    \hline
    \textbf{Kontroler} & \textbf{Q1 (Backlog)} & \textbf{Q2 (Resource)} & \textbf{Q3 (Balanced)} & \textbf{Regret} \\
    \hline
    Statik           & 240,1 & 833,6 & 851,9 & 2,477 \\
    \hline
    Rule-Based       & 147,2 & 374,4 & 412,3 & 0,785 \\
    \hline
    PID              & 95,8 & 268,1 & 296,8 & 0,180 \\
    \hline
    \textbf{LQR}     & \textbf{82,3} & \textbf{217,4} & \textbf{241,4} & \textbf{0,000} \\
    \hline
    \end{tabular}
\end{table}

LQR menghasilkan biaya terendah pada seluruh preset $Q$ dengan \emph{regret} 0{,}000.
PID berada di urutan kedua dengan \emph{regret} rata-rata 0{,}180, sementara Statik tertinggal
jauh dengan \emph{regret} 2{,}477.
Pada preset Q1\_backlog yang menempatkan prioritas tertinggi pada antrean, LQR (82,3) lebih
rendah dari PID (95,8) sebesar 14\% dan dari Statik (240,1) sebesar 66\%.

Pemecahan biaya per pola beban pada preset Q1 ditampilkan pada Gambar~\ref{fig:bar-pattern-full}.
Selisih antar kontroler paling kentara pada pola \emph{step} dan \emph{periodic\_step},
yaitu kondisi di mana antrean naik secara berkelanjutan sehingga kemampuan kontroler
menyesuaikan \emph{batch} sangat menentukan.

\begin{figure}[H]
    \centering
    \includegraphics[width=0.85\textwidth]{IV_R2__bar_cost_pattern_full_Q1.png}
    \caption{Fungsi biaya $J$ per pola beban pada preset Q1 (sumber daya penuh)}
    \label{fig:bar-pattern-full}
\end{figure}

Throughput per kontroler pada preset Q1 ditunjukkan pada Gambar~\ref{fig:throughput-full}.
LQR mencapai \emph{throughput} tertinggi pada pola \emph{impulse}, di mana antrean melonjak
hingga 10.000 dan kontroler perlu mempercepat drainase secara agresif.

\begin{figure}[H]
    \centering
    \includegraphics[width=0.85\textwidth]{IV_R3__bar_throughput_full_Q1.png}
    \caption{\emph{Throughput} per kontroler pada preset Q1 (sumber daya penuh)}
    \label{fig:throughput-full}
\end{figure}

Gambar~\ref{fig:timeseries-queue-full} menampilkan \emph{timeseries} antrean pada lima pola
beban dengan preset Q1.
LQR (warna hijau pada plot) konsisten menjaga antrean lebih rendah dibanding Statik dan
Rule-Based, terutama pada pola \emph{step} dan \emph{periodic\_step}.

\begin{figure}[H]
    \centering
    \includegraphics[width=\textwidth]{IV_R5__timeseries_queue_full_Q1.png}
    \caption{\emph{Timeseries} panjang antrean per kontroler pada lima pola beban (preset Q1, sumber daya penuh)}
    \label{fig:timeseries-queue-full}
\end{figure}

\subsection{Hasil Sumber Daya Setengah}
\label{sec:hasil-half}

Pada konfigurasi setengah, sistem mengonsumsi 1780 sampel dengan latensi rata-rata 2455~ms,
sekitar 1{,}3 kali latensi pada konfigurasi penuh.
Tabel~\ref{tab:cost-J-half} menunjukkan rangkuman fungsi biaya $J$ pada konfigurasi ini.

\begin{table}[H]
    \centering
    \caption{Nilai rata-rata fungsi biaya $J$ pada konfigurasi sumber daya setengah}
    \label{tab:cost-J-half}
    \renewcommand{\arraystretch}{1.4}
    \begin{tabular}{|l|r|r|r|r|}
    \hline
    \textbf{Kontroler} & \textbf{Q1 (Backlog)} & \textbf{Q2 (Resource)} & \textbf{Q3 (Balanced)} & \textbf{Regret} \\
    \hline
    Statik           & 157,8 & 483,7 & 513,3 & 1,418 \\
    \hline
    Rule-Based       & 154,2 & 386,5 & 432,9 & 1,101 \\
    \hline
    PID              & 87,0 & 226,1 & 254,7 & 0,150 \\
    \hline
    \textbf{LQR}     & \textbf{63,7} & \textbf{169,5} & \textbf{191,9} & \textbf{0,000} \\
    \hline
    \end{tabular}
\end{table}

Polanya sama dengan konfigurasi penuh, di mana LQR memberikan biaya terendah dan diikuti
oleh PID.
Yang menarik adalah jarak relatif antar kontroler.
Pada konfigurasi setengah, \emph{regret} Rule-Based menjadi 1{,}101 atau jauh lebih buruk
dibanding pada konfigurasi penuh yang hanya 0{,}785.
Ini menunjukkan bahwa Rule-Based semakin tertinggal saat sumber daya menipis, karena aturan
ambang batas yang statis tidak menyesuaikan diri terhadap kapasitas yang lebih kecil.

LQR justru tetap kompetitif dengan \emph{regret} 0{,}000 dan biaya absolut yang lebih rendah
dibanding versi penuh (63,7 vs 82,3 pada Q1).
Hal ini terjadi karena LQR menyesuaikan aksi kontrol berdasarkan kondisi sistem aktual,
sehingga ketika kapasitas berkurang, antrean terbangun lebih lambat dan kontroler memilih
batch yang lebih moderat tanpa kehilangan kemampuan drainase.

Gambar~\ref{fig:bar-pattern-half} menampilkan distribusi biaya per pola beban untuk
konfigurasi setengah, sedangkan Gambar~\ref{fig:timeseries-queue-half} menampilkan
\emph{timeseries} antrean.

\begin{figure}[H]
    \centering
    \includegraphics[width=0.85\textwidth]{IV_R7__bar_cost_pattern_half_Q1.png}
    \caption{Fungsi biaya $J$ per pola beban pada preset Q1 (sumber daya setengah)}
    \label{fig:bar-pattern-half}
\end{figure}

\begin{figure}[H]
    \centering
    \includegraphics[width=\textwidth]{IV_R8__timeseries_queue_half_Q1.png}
    \caption{\emph{Timeseries} panjang antrean per kontroler pada lima pola beban (preset Q1, sumber daya setengah)}
    \label{fig:timeseries-queue-half}
\end{figure}

Tabel~\ref{tab:robustness} merangkum perubahan \emph{regret} antara dua konfigurasi.
Statik dan Rule-Based degradasi cukup tajam saat sumber daya menipis, sementara PID dan LQR
relatif stabil.
Bahkan, \emph{regret} Statik membaik dari 2{,}477 ke 1{,}418 saat sumber daya berkurang.
Anomali ini muncul karena \emph{regret} dinormalisasi terhadap kontroler terbaik, sehingga
ketika baseline LQR juga melemah secara absolut, jarak relatif terhadap LQR menyempit.
Yang penting adalah peringkat tetap konsisten yaitu LQR di posisi terbaik diikuti PID.

\begin{table}[H]
    \centering
    \caption{Perbandingan \emph{regret} rata-rata pada dua konfigurasi sumber daya}
    \label{tab:robustness}
    \renewcommand{\arraystretch}{1.4}
    \begin{tabular}{|l|r|r|}
    \hline
    \textbf{Kontroler} & \textbf{Regret (Penuh)} & \textbf{Regret (Setengah)} \\
    \hline
    Statik       & 2,477 & 1,418 \\
    \hline
    Rule-Based   & 0,785 & 1,101 \\
    \hline
    PID          & 0,180 & 0,150 \\
    \hline
    \textbf{LQR} & \textbf{0,000} & \textbf{0,000} \\
    \hline
    \end{tabular}
\end{table}

\subsection{Adaptasi Antar Preset Q}
\label{sec:adaptasi}

Salah satu karakteristik LQR yang membedakannya dari kontroler lain adalah kemampuannya
mengubah perilaku berdasarkan matriks bobot $Q$.
Gambar~\ref{fig:state-across-q} menunjukkan rata-rata variabel sistem untuk seluruh kontroler
pada tiga preset $Q$ yang berbeda, dengan baseline ditampilkan sebagai garis putus-putus
karena nilainya konstan terhadap $Q$.

\begin{figure}[H]
    \centering
    \includegraphics[width=\textwidth]{IV_R4__state_across_q_full.png}
    \caption{Rata-rata variabel sistem per kontroler pada tiga preset $Q$ (sumber daya penuh)}
    \label{fig:state-across-q}
\end{figure}

Pada panel latensi, LQR konsisten berada di tingkat terendah dengan pergeseran tipis dari
1100 ms pada Q1 ke 1228 ms pada Q3.
Pada panel utilisasi CPU, LQR-Q2 yang memprioritaskan penghematan sumber daya memberikan
nilai paling rendah (22\%) dibanding LQR-Q1 dan LQR-Q3.
Pada panel I/O, LQR-Q2 juga menunjukkan operasi tulis paling sedikit, sejalan dengan
prioritas resource pada preset tersebut.

Kontroler Statik, Rule-Based, dan PID tampil sebagai garis horizontal karena perilakunya
tidak ditentukan oleh $Q$.
Ketiganya tidak memiliki mekanisme untuk menyerap perubahan prioritas, sehingga menghasilkan
nilai variabel sistem yang sama untuk seluruh preset.
Inilah yang menjadi nilai utama LQR sebagai kontroler berbasis fungsi biaya, yaitu satu
implementasi dapat dikonfigurasi ulang untuk prioritas yang berbeda hanya dengan mengganti
matriks $Q$.

Variabel utilisasi memori tidak ditampilkan karena pembacaan dari kontainer Elasticsearch
menghasilkan nilai nol pada eksperimen ini, sehingga tidak dapat digunakan untuk membandingkan
adaptasi antar kontroler.

\subsection{Robustness Lintas Mesin dan Konfigurasi}
\label{sec:robustness}

Untuk menguji apakah keunggulan LQR bertahan pada kombinasi perangkat keras dan kapasitas
yang berbeda, eksperimen yang sama dijalankan pada empat kombinasi (dua mesin uji, masing-masing
dengan dua konfigurasi sumber daya).
Mesin pertama menggunakan AMD Ryzen sebagai host, sementara mesin kedua menggunakan Apple
Silicon.
Pada setiap mesin, Elasticsearch dijalankan pada konfigurasi penuh (2{,}0 CPU, 2 GiB memori,
heap 1 GiB) dan konfigurasi setengah (1{,}0 CPU, 1 GiB memori, heap 512 MiB).

Tabel~\ref{tab:cross-experiment} merangkum nilai $J$ rata-rata setiap kontroler pada keempat
eksperimen.
Nilai dihitung dengan merata-ratakan cost $J$ atas lima pola beban dan tiga preset $Q$ untuk
setiap kontroler, dengan varian LQR dipasangkan pada preset $Q$ yang sesuai.

\begin{table}[H]
    \centering
    \caption{Cost $J$ rata-rata per kontroler pada empat eksperimen}
    \label{tab:cross-experiment}
    \renewcommand{\arraystretch}{1.4}
    \begin{tabular}{|l|r|r|r|r|}
    \hline
    \textbf{Kontroler} & \textbf{Mac Full} & \textbf{Mac Half} & \textbf{Ryzen Full} & \textbf{Ryzen Half} \\
    \hline
    Statik           & 962,4 & 4697,7 & 640,7 & 385,0 \\
    \hline
    Rule-Based       & 237,4 & 2805,8 & 317,2 & 324,4 \\
    \hline
    PID              & 63,2 & 3575,7 & 209,7 & 175,9 \\
    \hline
    \textbf{LQR}     & \textbf{16,2} & \textbf{1254,2} & \textbf{177,9} & \textbf{159,6} \\
    \hline
    \end{tabular}
\end{table}

LQR memberikan cost $J$ terendah di seluruh empat eksperimen, sedangkan Statik selalu
menjadi yang tertinggi.
Pada Mac Half, biaya seluruh kontroler melonjak jauh lebih tinggi dibanding eksperimen lain,
yang menunjukkan bahwa kombinasi perangkat keras Mac dengan kapasitas terbatas menyulitkan
seluruh strategi untuk menjaga deviasi state.
Bahkan dalam kondisi sulit tersebut, LQR (1254,2) tetap setengah dari Statik (4697,7)
dan jauh di bawah PID (3575,7).

Gambar~\ref{fig:cross-experiment} menampilkan visualisasi yang sama dalam bentuk batang.
Setiap panel menggunakan skala mandiri karena nilai absolut $J$ berbeda jauh antar eksperimen,
sehingga skala bersama akan membuat perbedaan dalam panel kecil tidak terbaca.

\begin{figure}[H]
    \centering
    \includegraphics[width=\textwidth]{IV_R9__cross_experiment_cost.png}
    \caption[Cost J kontroler lintas empat eksperimen]{Cost $J$ rata-rata per kontroler pada empat eksperimen, dengan setiap panel berskala independen}
    \label{fig:cross-experiment}
\end{figure}

Perlu dicatat bahwa nilai $J$ absolut tidak dapat dibandingkan secara langsung antar mesin
karena matriks $Q$ untuk Mac dan Ryzen dikalibrasi terpisah dari hasil identifikasi sistem
masing-masing.
Yang konsisten adalah ranking dalam setiap eksperimen, di mana LQR selalu di urutan teratas.
Hal ini menunjukkan bahwa keunggulan LQR bukan artefak dari satu konfigurasi tertentu,
melainkan akibat langsung dari mekanisme pengendalian yang menggabungkan seluruh variabel
sistem ke dalam satu fungsi biaya.

\subsection{Kasus Batas pada Sumber Daya Quarter}
\label{sec:quarter}

Konfigurasi seperempat (\emph{Quarter}) dengan 0{,}5 CPU, 512 MiB memori, dan JVM heap 256 MB
diuji sebagai kasus batas untuk melihat perilaku sistem ketika kapasitas Elasticsearch tidak
lagi memadai.
Latensi rata-rata pada konfigurasi ini mencapai 6797~ms atau sekitar 3{,}7 kali latensi pada
konfigurasi penuh, yang menunjukkan bahwa Elasticsearch sudah dalam kondisi saturasi.
Pada kondisi ini, perbedaan antar kontroler menjadi kurang relevan karena bottleneck bergeser
dari pengaturan parameter ke kapasitas pemrosesan basis data itu sendiri.
Hasil pada konfigurasi Quarter tidak dilampirkan dalam bentuk tabel atau plot karena
karakteristiknya sebagai kasus degradasi yang tidak mencerminkan kondisi operasional yang ditargetkan.

